\chapter{Parity and Permanent Deviations}
\label{ch:parity-philosophy}

Every project this book has covered draws a line between what it will faithfully reproduce
and what it will not. Sharp-runtime states its target as \emph{maximum practical parity}:
public API, semantics, defaults, error messages, and algorithms should match .NET as closely
as C\texttt{++} allows. That target is stronger than name compatibility and narrower than
reimplementing the CLR. Its permanent deviations are part of the contract, not embarrassing
footnotes.

\section{Three outcomes, not one generic stub}

An implementable operation should behave like .NET. An operation whose mechanism cannot be
represented should throw \cnaclass{NotImplementedException} with an explanation rather than
return a plausible but wrong value. A meaningful operation that the current host cannot
provide should instead throw \cnaclass{PlatformNotSupportedException}. These outcomes answer
different questions:

\begin{description}
  \item[Parity] the operation is implemented and its observable behavior is the verification
  target;
  \item[Permanent deviation] the CLR-dependent capability is outside the runtime's design;
  \item[Platform limitation] the abstraction is valid, but this platform lacks its mechanism.
\end{description}

A silent success value is not a fourth acceptable category. This taxonomy lets callers and
tests distinguish a design boundary from an environment boundary.

\section{Reflection: correctly a stub, not a gap}

\cnaclass{System.Type}, \cnaclass{Activator}, and enum reflection (\texttt{Enum.GetNames} /
\texttt{GetValues}) are, in the project's own words, ``completely out of scope. Stubs are the
correct end state.'' This is a stronger claim than ``not yet implemented'' --- it states that
a stub \emph{is} the intended final behavior, not a placeholder awaiting more engineering.
Reflection depends on runtime type metadata a compiled C\texttt{++} binary simply does not
carry the way a .NET assembly does; reproducing it would mean building a real, if partial,
runtime type system from scratch, disproportionate to what CNA's own XNA-facing code actually
needs.

\section{Garbage collection: inert compatibility answers, by design}

\cnaclass{System.GC} has no collector behind it. Mutating methods such as \texttt{Collect} and
\texttt{SuppressFinalize} are no-ops, while queries return fixed documented answers: allocation
statistics are zero, no-GC-region requests return false, notification waits report
\texttt{NotApplicable}, and \texttt{MaxGeneration} remains the compatibility constant 2. Object
lifetime is managed by RAII and \texttt{shared\_ptr}; callers may compile defensive GC calls but
must not infer live heap telemetry from the surface.

\section{Delegates: a three-tier model, corrected once in the project's own history}

This is worth covering in detail because the project's own documentation records a real
self-correction: an earlier description of the delegate model was found inaccurate during an
audit and rewritten. The corrected, current model has three distinct tiers, each solving a
different shape of problem. Most delegate-shaped C\# types --- \texttt{Action},
\texttt{Func}, most \texttt{*Callback} and \texttt{*EventHandler} aliases --- are bare
\texttt{using} aliases over
\texttt{std::function}: single-target, no multicast, no \texttt{BeginInvoke} /
\texttt{EndInvoke}, sufficient for the overwhelming majority of call sites that only ever
needed one callback anyway. \cnaclass{System::Delegate} is a real multicast base
(\texttt{Combine}, \texttt{Remove}, \texttt{GetInvocationList}) for the smaller set of call
sites that genuinely need C\#'s multicast-delegate combining semantics.
\cnaclass{System::MulticastAction<Args...>} (Chapter~\ref{ch:sharp-runtime-overview}) is a
third, purpose-built type specifically for multicast event \emph{fields}, using
token-based removal since C\texttt{++} callables have no target-plus-method equality to
compare by the way C\# delegates do. The specialized
\texttt{EventHandler<TEventArgs>} collection applies the same event-field idea to the
sender/argument shape and adds an optional replay hook. Of these representations, only
\cnaclass{Delegate} exposes \texttt{DynamicInvoke}; it always throws
\texttt{NotImplementedException}. The alias and event-field types do not pretend to expose a
late-bound argument-array call at all.

\begin{lstlisting}
// Worked example: all three tiers side by side, each solving the shape of
// problem it exists for.

// Tier 1 -- Action/Func alias: single-target, the overwhelming common case.
System::ActionT<int> onScoreChanged = [](int newScore) { UpdateHudScore(newScore); };
onScoreChanged(100);

// Tier 2 -- System::Delegate: real multicast combining semantics, when a
// caller genuinely needs to build up and later shrink an invocation list.
auto d1 = std::make_shared<System::Delegate>(System::Delegate::ErasedInvoke([] { PlayFanfare(); }));
auto d2 = std::make_shared<System::Delegate>(System::Delegate::ErasedInvoke([] { UnlockNextLevel(); }));
std::shared_ptr<System::Delegate> onLevelComplete = System::Delegate::Combine(d1, d2);
onLevelComplete->Invoke();                                    // runs both
onLevelComplete = System::Delegate::Remove(onLevelComplete, d1);  // now only d2 remains

// Tier 3 -- MulticastAction: the right choice for an event FIELD, since
// std::function values have no target-plus-method equality to remove by.
System::MulticastAction<> onGameOver;
auto token = onGameOver.Add([] { ShowGameOverScreen(); });
onGameOver.Remove(token);
\end{lstlisting}

\section{Serialization, P/Invoke, and cryptography: three different reasons for the same conclusion}

Serialization is ``ignored, not needed for game code'' --- a scope judgment rather than a
technical impossibility. P/Invoke and general native interop are out of scope for the same
underlying reason Chapter~\ref{ch:design-philosophy}'s type-alias table exists in the first
place: this project already \emph{is} native code, so a marshaling layer for calling into
native code from managed code has no problem left to solve here.

Cryptography and TLS are the one deviation with the most fully spelled-out reasoning of all,
and it deserves quoting closely: symmetric and asymmetric cryptography, X.509 certificates,
and \cnaclass{SslStream} are out of scope by an explicit, dated project-owner decision, on the
grounds that ``implementing this correctly needs either a large new external dependency
(OpenSSL/mbedTLS) or a hand-rolled, security-critical implementation, neither of which is
worth it for game code.'' This is a materially different kind of reasoning than the other
deviations in this chapter --- not ``this feature has no meaning here'' (reflection, GC) and
not ``this is out of scope for capacity reasons'' (serialization), but an explicit refusal to
carry the security liability of a hand-rolled cryptographic implementation, or the dependency
weight of a real one, for a class of application that does not need transport security built
into its own runtime layer. Digest and derivation APIs --- MD5, SHA, HMAC, and PBKDF2 ---
remain in scope. That does not make them ``non-cryptographic,'' nor does it imply that HMAC or
PBKDF2 has no key or password material. The narrower project distinction is that these APIs
do not implement reversible confidentiality, certificate validation, or a TLS channel.
Callers still own algorithm choice, parameter strength, and security review.

\section{IDisposable and the exception hierarchy: parity by pattern, not by mechanism}

Two further deviations are less often stated outright as deviations, because the project
does not treat them as ones --- they are cases where matching .NET's \emph{observable
behavior} required deliberately \emph{not} matching its mechanism. \cnaclass{IDisposable}'s
own header comment states the tradeoff directly: ``In C\#, the \texttt{using} statement calls
\texttt{Dispose()} automatically at the end of the block. In C\texttt{++}, use RAII or call
\texttt{Dispose()} explicitly.'' There is no C\texttt{++} language construct that runs
arbitrary code at scope exit the way a \texttt{using} block does short of writing a dedicated
RAII wrapper per call site, so this project's actual parity target is narrower and more
precise than ``reproduce the \texttt{using} statement'': reproduce \emph{Dispose()'s own
contract} --- idempotent, generally non-throwing, and safe to call more than once --- and let the
caller choose their own C\texttt{++}-native mechanism (an explicit call, a destructor, or a
scope guard) for invoking it at the right time.

\cnaclass{System::Threading::ReaderWriterLockSlim} is a concrete, real implementation of this
contract worth reading closely, because its own header comment documents a genuine
self-correction of exactly the kind this chapter's delegate discussion above already showed
this project performing on itself. An earlier version of this
class had three real, independently-found bugs, each one a case where a C\texttt{++}
implementation that merely \emph{looked} equivalent to the real .NET semantics was not
actually equivalent under closer verification: it discarded the
\texttt{millisecondsTimeout} parameter on every \texttt{TryEnter*} overload, always making a
single non-blocking attempt regardless of what the caller passed; it ignored the constructor's
\texttt{LockRecursionPolicy} entirely, so same-thread recursive acquisition always deadlocked
instead of throwing \texttt{LockRecursionException} under the real .NET-default
\texttt{NoRecursion} policy; and it tracked reader-lock ownership by set \emph{membership}
rather than by count, so a legitimately nested
\texttt{EnterReadLock()} / \texttt{EnterReadLock()} / \texttt{ExitReadLock()} sequence's second
\texttt{ExitReadLock()} call threw \texttt{SynchronizationLockException} instead of correctly
decrementing --- which additionally meant the internal reader tally could never reach zero
from a genuinely nested acquisition, permanently starving any waiting writer. All three are
fixed in the current header, verified directly against real .NET's own
\texttt{TryEnterReadLockCore} / \texttt{TryEnterWriteLockCore} / \texttt{TryEnterUpgradeableReadLockCore}
logic per the header's own verification comment.

\begin{lstlisting}
// Worked example: IDisposable's contract, applied explicitly (no C++ "using"
// block exists), with the real ObjectDisposedException guard every public
// entry point in ReaderWriterLockSlim shares (throwIfDisposed(), called
// first inside TryEnterReadLock/TryEnterWriteLock/TryEnterUpgradeableReadLock).
System::Threading::ReaderWriterLockSlim rwLock;

rwLock.EnterReadLock();
// ... read shared state ...
rwLock.ExitReadLock();

rwLock.EnterWriteLock();
// ... mutate shared state ...
rwLock.ExitWriteLock();

rwLock.Dispose();          // explicit -- no using-block equivalent exists
rwLock.Dispose();          // safe to repeat: Dispose() only ever sets a flag
try {
    rwLock.EnterReadLock();      // throwIfDisposed() now fires
} catch (const System::ObjectDisposedException&) {
    // "ReaderWriterLockSlim" -- the disposed object's own class name,
    // matching real .NET's ObjectDisposedException(objectName) constructor.
}
\end{lstlisting}

The exception hierarchy this last block leans on is itself worth a concrete scale rather than a
vague gesture at ``lots of exception types'': a project-wide search finds roughly 111
exception classes across the \texttt{include/System} tree, the overwhelming majority deriving
not directly from \cnaclass{System::Exception} but through the intermediate
\cnaclass{System::SystemException} base --- matching real .NET's own two-level shape, where
framework-thrown exceptions (\texttt{ArgumentException}, \texttt{InvalidOperationException},
\texttt{ObjectDisposedException}, and the rest this book's other chapters have already used by
name) sit under \texttt{SystemException} rather than hanging directly off the root. The root
class's own header comment makes explicit exactly which real .NET \texttt{Exception} members
this project deliberately does not provide, and why each one is a concrete instance of a
deviation this chapter has already stated abstractly: \texttt{GetBaseException()} would
require cloning the thrown object through its base type, which needs a virtual-clone mechanism
this port does not have; \texttt{ToString()} would need a \texttt{GetType()}-equivalent class
name, out of scope under this chapter's opening reflection deviation; and \texttt{TargetSite}
and \texttt{GetObjectData} are out of scope for the same reflection and serialization
deviations covered earlier in this chapter. Nothing here is a new decision --- it is the same
two decisions, now traceable to the one root class every exception in the ecosystem inherits
from.

\section{Naming is a compatibility surface}

Sharp-runtime uses PascalCase .NET names, but C\texttt{++} cannot expose C\# properties
directly. The binding convention is \texttt{getXxxProperty()} and
\texttt{setXxxProperty()}; indexers are the systematic exception and use
\texttt{getItem()} / \texttt{setItem()}. Public headers contain thousands of these accessors,
so a cosmetic refactor would be an ecosystem break rather than a local cleanup.

The project's own rules name CNA as the reason broad header renaming is frozen. That is strong
evidence of co-evolution: the runtime is a separate repository, yet its source conventions
are constrained by CNA's compiled public surface. Porting rules similarly require the
\texttt{intcs}/\texttt{bytecs} family and direct nested namespace syntax, while LINQ-shaped
work should use \texttt{std::ranges} rather than growing a second query runtime.

When a naming correction is nevertheless intentional, the project does not preserve a stale
alias indefinitely. The \texttt{getCurrent()} to \texttt{getCurrentProperty()} migration
changed 61 call sites across 23 files in one commit and required downstream repair. The lesson
is not that names may churn casually; it is that a deliberate convention repair should be
atomic, auditable, and free of permanent compatibility debris.

\section{Two further, narrower accepted boundaries}

Platform behavior is split more finely than a single ``POSIX-only'' label. Sockets and
\cnaclass{IO::RandomAccess} have Windows and POSIX implementations and throw on Emscripten.
\texttt{AppDomain}'s base directory works on Windows, macOS and Linux/POSIX, with a relative
virtual-filesystem fallback on Emscripten; \texttt{TimeZoneInfo} likewise has Windows and POSIX
paths but only UTC/local fallback there. \texttt{Diagnostics::Process} and POSIX-signal
registration are POSIX implementations that throw on Windows/Emscripten, while
\texttt{NetworkInterface} enumeration and \texttt{FileSystemWatcher} are narrower Linux-only
operations. Unsupported paths compile and fail explicitly rather than disappearing from the
headers.

Separately, \texttt{Decimal}, \texttt{Int128},
and \texttt{UInt128} require a compiler-provided 16-byte \texttt{\_\_int128}. The CMake probe
admits x86-64 GCC/Clang, including x86-64 MinGW GCC, but excludes MSVC and i686 MinGW GCC. This
is a compiler-capability boundary, not simply a Windows boundary, and reflects the deliberate,
dated decision not to hand-roll 128-bit arithmetic for toolchains that lack the native type.

\section{A documented deviation still needs verification}

Calling a gap permanent does not exempt its boundary from tests. \cnaclass{GC}'s inert methods
and sentinel queries should remain callable and stable; \cnaclass{Type}'s identity operations
must work even though its classification
predicates are placeholders; \cnaclass{Delegate::DynamicInvoke} must fail explicitly; and
unsupported platform operations must throw the promised exception rather than silently
degrade. Chapter~\ref{ch:sharp-verification-audit} shows how component tests, negative
consumers, sanitizer runs, and the repository's large audit ledger test both implemented
behavior and the honesty of these boundaries.
